\documentclass[10pt,a4paper]{amsart}
\usepackage[margin=19mm]{geometry}
\usepackage{amsmath,amssymb,mathtools}
\usepackage[scr=rsfs,cal=euler]{mathalfa}
\usepackage{xfrac}
\usepackage[unicode,hidelinks,bookmarks=false]{hyperref}
\newcommand{\ie}{i.e.\ }
\newcommand{\cf}{cf.\ }
\newcommand{\resp}{respectively\ }
\newcommand{\Subsection}{\subsection}
\providecommand{\nobreakdash}{\nobreak\hbox{-}}
\setlength{\emergencystretch}{2em}
\multlinegap0pt
\usepackage{bm}
\usepackage{bbm}
\usepackage{dsfont}
\usepackage{xspace}
\usepackage{cancel}
\usepackage{mathtools}
\usepackage{amsthm}
\makeatletter
\@ifundefined{theorem}{\newtheorem{theorem}{Theorem}}{}
\makeatother
\title[Mapping Mathematical Hardness: Machine-Assisted Conjecture D]{Mapping Mathematical Hardness: Machine-Assisted Conjecture Discovery and the Quantification of Non-Triviality}
\author{Madhuparna Das}
\date{Source: arXiv:2606.14804v2 (CC BY 4.0)}
\begin{document}
\maketitle

\noindent\emph{Source and licence.} Mapping Mathematical Hardness: Machine-Assisted Conjecture Discovery and the Quantification of Non-Triviality. Version arXiv:2606.14804v2 (CC BY 4.0), distributed under the Creative Commons Attribution 4.0 International licence, as recorded in the arXiv licence metadata for this exact version and shown on the abstract page https://arxiv.org/abs/2606.14804v2. Full rights notice: licenses/arxiv-2606.14804-CC-BY-4.0.txt.\par\smallskip

\noindent\textit{Editorial scope.} Counted unit: the Theorem whose source label is \texttt{primedrop} in the article (source0.tex lines 343--348, source label \texttt{primedrop}), delivered together with the article's own complete original proof of it (source0.tex lines 350--414, one \texttt{proof} environment of 65 lines). No mathematics is written, repaired or re-derived: statement and proof are the article's own text line for line. Required context retained uncounted: monotonerefinement (lines 279--284). Every definition, display and label of those lines is retained unchanged. Omitted from this package and not redistributed: 1--278, 285--342, 349--349, 415--1227. Nothing inside a retained excerpt is omitted or altered: every retained body is delivered exactly as the article's lines are written, so no omission receipt of any kind accompanies this package. Citations: the retained excerpts carry no citation command at all, so no bibliography entry of the article is transcribed and no reference is redistributed. Reference adaptation: none was needed and none was made; every reference inside the delivered unit resolves inside it, so no unresolved reference or citation is produced. Printed slips kept literal: no printed word, symbol, hypothesis or number of the counted statement or of its proof was altered. Layout and class: the article's own class is \texttt{amsart} with packages including graphicx, amsthm, amsmath, amssymb, amsfonts, mathtools, fullpage, python, enumitem, bm, tikz, algpseudocode, algorithm2e, appendix; the wrapper is the lane's pinned \texttt{amsart} editorial wrapper at 10pt on A4, which restates the theorem-like environments the retained text needs and adds the article's own macro definitions that the retained text needs (none), with extra wrapper packages (bm, bbm, dsfont, xspace, cancel, mathtools). The delivered numbering is the unit's own. Assets: no retained span contains a figure environment, a graphics-inclusion command, a PDF-page inclusion, a table or a TikZ picture; no image file is copied into this package, the package declares no asset dependency and no image file is redistributed with it. Licence: version arXiv:2606.14804v2 of this article is distributed under the Creative Commons Attribution 4.0 International licence, as recorded in the arXiv licence metadata for this exact version and shown on the abstract page https://arxiv.org/abs/2606.14804v2. A full rights notice is delivered at licenses/arxiv-2606.14804-CC-BY-4.0.txt. Characters: every retained excerpt is pure ASCII, so the delivered engine, XeLaTeX with its default Latin Modern text font, reports no missing character.
\begin{theorem}[HypothesiX, Conjecture A.3]\label{monotonerefinement}
If $Q_1 \mid Q_2$ are both square-free and divisible by 6, then for every $x \geq 7$
\begin{align*}
B_{Q_2}(x) \leq B_{Q_1}(x).    
\end{align*}
\end{theorem}

\begin{theorem}[HypothesiX, Conjecture A.8]\label{primedrop}
Fix squarefree $Q$ and prime $p \geq 5$ coprime to $Q$. There exists $X_0$ such that for all $x \geq X_0$
\begin{align*}
B_{pQ}(x)\leq B_Q(x)-1    
\end{align*}          
\end{theorem} 

\begin{proof}
Choose any $s_o\in U_Q$. Then by the Chinese remainder theorem each $s\in U_Q$, lifts exactly $p$ distinct classes modulo $pQ$. The fiber over $s_0$ is
\begin{align*}
v_{s_0}
=\{r\in U_{pQ}:r\equiv s_0(\bmod{Q})\},
\end{align*}
consisting of the $p-2$ lifts satisfying $r\not\equiv0,-2(\bmod{p})$. Now, we define
\begin{align*}
r_0
\coloneqq F(s_0,Q,p-2,p),
\end{align*}
where $F(s_0,Q,p-2,p)$ means the unique residue class $r_0(\bmod{pQ})$ such that
\begin{align*}
r_0\equiv s_0(\bmod{Q})\quad\text{and}\quad r_0\equiv p-2(\bmod{p}).    
\end{align*}
Similarly, 
\begin{align*}
r_3
\coloneqq F(s_0+2,Q;2,p).
\end{align*}

\noindent
Note that $p\geq5$ and $(p,Q)=1$. Then, for $r_0\equiv p-2(\bmod{p})$, $(r_0,p)=1$ and $(r_0,Q)=(s_0,Q)=1$. Analogously, $(r_3,p)=(r_3,Q)=(s_0+2,Q)=1$. By Dirichlet's theorem, we can write
\begin{align*}
\pi(p,Q;r_0)\geq 1\quad\text{and}\quad
\pi(p,Q;r_3)\geq 1.
\end{align*}

\noindent
Since $r_0\equiv p-2\equiv -2(\bmod{p})$, we have $p|(r_0+2)$, so $r_0\notin U_{pQ}$ and thus $r_0\notin v_{s_0}$. Therefore, $r_0$ contributes to $\pi(x;Q,s_0)$ but not to $\sum_{r\in v_{s_0}}\pi(x;pQ,r)$. For any $r\in v_{s_0}$, we have $r \not\equiv -2(\bmod{p})$ and $r\not\equiv 0(\bmod{p})$, 
so $r+2\not\equiv 2(\bmod{p})$. But $r_3\equiv 2(\bmod{p})$, so $r_3$ is not of the form $r+2$ for any $r\in v_{s_0}$. Since $r_3 \equiv s_0+2(\bmod{Q})$, it does contribute to $\pi(x,;Q,s_0+2)$ but not to $\sum_{r\in v_{s_0}}\pi(x;pQ,r+2)$. 

\medskip
Set
\begin{align*}
A
\coloneqq\sum_{r\in v_{s_0}}\pi(x;pQ,r)\quad\text{and}\quad
B
\coloneqq\sum_{r\in v_{s_0}}\pi(x;pQ,r+2).
\end{align*}
Then for $x\geq X_0$, we have
\begin{align*}
\pi(x;Q,s_0)
\geq A +\pi(x;pQ,r_0)\geq A+1.
\end{align*}
Similarly, $\pi(x;Q,s_0+2)\geq B+1$. Since the $\min$ function is monotone increasing
\begin{align*}
\min\bigl(\pi(x; Q, s_0),\pi(x; Q, s_0 + 2)\bigr)
\geq\min(A+1, B+1) 
=\min(A,B)+1.    
\end{align*}
By the inequality $\sum_{r}\min(a_r,b_r)\leq\min(\sum_{r}a_r,\sum_{r}b_r)$ for $a_r,b_r\geq 1$, yields
\begin{align*}
\min\bigl(\pi(x;Q,s_0),\pi(x; Q, s_0 + 2)\bigr)
\geq\sum_{r\in v_{s_0}}\min\bigl(\pi(x;pQ,r),\pi(x;pQ,r+2)\bigr)+1.    
\end{align*}
Summing over all $s\in U_Q$ and using non-negativity of each term from Theorem~\ref{monotonerefinement}, 
\begin{align*}
B_Q(x)-B_{pQ}(x)
=&\sum_{s\in U_Q}\min\bigl(\pi(x;Q,s),\pi(x;Q,s+2)\bigr)\\
&-\sum_{s\in U_Q}\sum_{r\in v_s} \min\bigl(\pi(x;pQ,r),\pi(x;pQ,r+2)\bigr)
\ge 1,    
\end{align*}
thereby concluding the proof.
\end{proof}

\end{document}
